% ═══════════════════════════════════════════════════════════════════════
\section{Pourquoi \texorpdfstring{\LaTeX}{LaTeX} ?}
% ═══════════════════════════════════════════════════════════════════════

\begin{frame}{Le standard de l'écriture mathématique}
  \small
  \begin{columns}[c]
    \begin{column}{0.5\textwidth}
      \begin{itemize}
        \item Créé par Donald Knuth (\TeX, 1978) puis Leslie Lamport
          (\LaTeX, 1984), pour composer des mathématiques avec la qualité
          d'un livre imprimé.
        \item Utilisé pour presque tous les articles, manuels et notes de
          cours de mathématiques --- y compris les notes de ce cours.
        \item Gratuit, ouvert, et le même fichier compile partout, dans
          dix ans comme aujourd'hui.
      \end{itemize}
    \end{column}
    \begin{column}{0.46\textwidth}
      \begin{tcolorbox}[codestyle]
        \centering
        $\displaystyle \sum_{k=1}^{n} k^2 = \frac{n(n+1)(2n+1)}{6}$
        \\[8pt]
        $\displaystyle \int_0^{\pi} \sin x \,\mathrm{d}x = 2$
        \\[8pt]
        $\displaystyle \lim_{n\to\infty}\left(1+\frac1n\right)^{\!n} = \mathrm{e}$
      \end{tcolorbox}
      \centering\scriptsize Trois lignes de code chacune, sans souris.
    \end{column}
  \end{columns}
\end{frame}

\begin{frame}{Traitement de texte ou \LaTeX ?}
  \small
  \begin{center}
  \begin{tabular}{p{3.2cm}p{4.3cm}p{4.3cm}}
    \toprule
    & Traitement de texte & \LaTeX \\
    \midrule
    Formules & éditeur d'équations, à la souris & tapées au clavier, rendu
      de qualité livre \\
    Mise en page & vous la faites, à la main & calculée pour vous,
      cohérente partout \\
    Numérotation & à corriger quand on insère & automatique (sections,
      questions, figures) \\
    Format du fichier & binaire & texte brut : lisible, réutilisable,
      généré par l'IA \\
    Ce qu'on voit & le résultat, en direct & le code, puis le PDF
      compilé \\
    \bottomrule
  \end{tabular}
  \end{center}
  \vspace{4pt}
  On ne dessine pas la page : on décrit la structure (un titre, une
  question, une formule) et \LaTeX\ s'occupe de la mise en page.
\end{frame}

\begin{frame}{Pourquoi c'est particulièrement utile en enseignement}
  \begin{Classe}{Ce que \LaTeX\ change pour un enseignant de mathématiques}
    \small
    \begin{itemize}
      \item Un examen et son corrigé dans un seul fichier : un
        interrupteur affiche ou cache les solutions.
      \item Le barème se calcule tout seul : points par question,
        total, grille de correction.
      \item Des banques de questions réutilisables d'une année à l'autre,
        d'un groupe à l'autre (version A, version B).
      \item Des figures précises (TikZ) : triangles, repères, graphes de
        fonctions, à la bonne échelle.
      \item Des documents cohérents pour toute une équipe (école ou département) :
        mêmes en-têtes, même notation, même mise en page.
      \item Avec l'IA, le premier jet ne prend que quelques minutes ;
        le travail d'enseignant, c'est d'en vérifier et d'en ajuster le
        contenu.
    \end{itemize}
  \end{Classe}
\end{frame}

\begin{frame}{Ce qu'il faut accepter}
  \small
  \begin{itemize}
    \item \LaTeX\ n'est pas un logiciel « ce que vous voyez est ce que vous
      obtenez » : on écrit du code, on compile, on regarde le PDF.
    \item Une petite faute de frappe (une accolade oubliée) peut empêcher
      la compilation.
    \item L'IA écrit très bien du \LaTeX, mais elle se trompe aussi :
      commandes inventées, extensions inutiles, mathématiques fausses,
      texte en anglais.
  \end{itemize}
  \vspace{4pt}
  \begin{Rem}{L'objectif d'aujourd'hui}
    Pas de mémoriser des centaines de commandes. Comprendre assez la
    structure d'un fichier pour lire ce que l'IA produit, le modifier à la
    main sans tout casser, et réparer une erreur de compilation.
  \end{Rem}
\end{frame}
